% conclusion.tex -- Cross-Enterprise Verification Paper (RU-V7)

\section{Conclusion}

This paper presented a privacy-preserving cross-enterprise verification
protocol for ZK validium systems. The protocol uses a cross-reference
circuit---68,868 R1CS constraints performing dual Merkle path verification
with a Poseidon interaction commitment---to prove the consistency of
transactions recorded in two independent enterprise sparse Merkle trees
without revealing any data beyond the interaction's existence.

Our contributions are:

\begin{enumerate}
\item \textbf{Cross-reference circuit design}: A Groth16 circuit with 3
  public inputs and dual depth-32 Merkle path verification that binds
  two enterprise state roots through a hiding commitment. Proof generation
  takes 448~ms with rapidsnark.

\item \textbf{Three verification approaches}: Sequential (1.41$\times$
  overhead, no coordination), Batched Pairing (0.64$\times$, hub
  coordinator required), and Hub Aggregation (1.16$\times$, SnarkPack/Nebra
  UPA). All achieve $< 2\times$ overhead for the primary enterprise
  scenario (2--50 enterprises, linear interactions).

\item \textbf{Privacy analysis}: The protocol achieves 1 bit leakage per
  interaction (existence only), which is the information-theoretic minimum
  for any cross-enterprise verification on a public ledger. Four
  adversarial privacy tests confirm that no data content is exposed through
  the interaction commitment.

\item \textbf{Scalability characterization}: Batched pairing maintains
  sub-1$\times$ overhead up to 50 enterprises. Sequential verification
  stays below 2$\times$ for sparse interactions but exceeds the target in
  dense interaction graphs, establishing clear deployment guidelines.
\end{enumerate}

\subsection{Future Work}

\textbf{Multi-party interactions}: Extending the cross-reference circuit to
$k$-party interactions ($k > 2$) where a single proof verifies consistency
across $k$ enterprise trees simultaneously.

\textbf{Temporal privacy}: Mitigating interaction frequency leakage through
batched cross-reference submission with dummy proofs, preventing adversaries
from inferring the temporal pattern of inter-enterprise transactions.

\textbf{PLONK migration}: Transitioning to PLONK-based proving (Halo2 or
Goblin Plonk) to enable heterogeneous proof aggregation via StarkPack and
native recursive composition, eliminating the per-circuit trusted setup
requirement.

\textbf{On-chain validation}: Compiling the cross-reference circuit with
Circom, generating actual Groth16 proofs, and measuring on-chain verification
gas on Subnet-EVM to validate the analytical model.

\textbf{Formal verification}: The downstream pipeline (Logicist, Architect,
Prover) will formalize the cross-enterprise verification protocol in TLA+,
implement it in production TypeScript, and prove isolation and consistency
properties in Coq.
